\chapterstar{Acrónimos}

\begin{longtable}{|p{2.5cm}|p{11cm}|}

\hline
\rowcolor{gray!20}
\textbf{Acrónimo} & \textbf{Significado} \\
\hline
\endfirsthead

\hline
\rowcolor{gray!20}
\textbf{Acrónimo} & \textbf{Significado} \\
\hline
\endhead

\hline
\endlastfoot

AP    & Access Point — Punto de acceso inalámbrico que transmite la señal Wi-Fi a los dispositivos en una red local. \\ \hline
API   & Application Programming Interface — Interfaz de Programación de Aplicaciones. Conjunto de reglas y rutas que permiten la comunicación estandarizada entre sistemas de software. \\ \hline
BLE   & Bluetooth Low Energy — Bluetooth de Baja Energía. Estándar de comunicación inalámbrica de corto alcance y bajo consumo energético. \\ \hline
CDN   & Content Delivery Network — Red de Distribución de Contenido. Infraestructura de servidores distribuidos geográficamente que reduce los tiempos de carga de aplicaciones web. \\ \hline
CSS   & Cascading Style Sheets — Hojas de Estilo en Cascada. Lenguaje que define la presentación visual de documentos HTML. \\ \hline
ERP   & Enterprise Resource Planning — Planificación de Recursos Empresariales. Sistema integrado de gestión que centraliza los procesos de una organización. \\ \hline
ESIME & Escuela Superior de Ingeniería Mecánica y Eléctrica. Unidad académica del Instituto Politécnico Nacional. \\ \hline
GPIO  & General Purpose Input/Output — Entrada/Salida de Propósito General. Pines digitales de un microcontrolador que pueden configurarse como entradas o salidas. \\ \hline
HTML  & HyperText Markup Language — Lenguaje de Marcado de Hipertexto. Lenguaje estándar para estructurar contenido en páginas web. \\ \hline
HTTP  & HyperText Transfer Protocol — Protocolo de Transferencia de Hipertexto. Protocolo de comunicación base de la World Wide Web. \\ \hline
HTTPS & HyperText Transfer Protocol Secure — Protocolo Seguro de Transferencia de Hipertexto. Versión cifrada del protocolo HTTP mediante TLS/SSL. \\ \hline
I²C   & Inter-Integrated Circuit — Circuito Inter-Integrado. Protocolo de comunicación digital de corto alcance que permite conectar múltiples dispositivos mediante dos líneas: datos (SDA) y reloj (SCL). \\ \hline
ICE   & Ingeniería en Comunicaciones y Electrónica. Programa académico de la ESIME Culhuacán del Instituto Politécnico Nacional. \\ \hline
IDE   & Integrated Development Environment — Entorno de Desarrollo Integrado. Aplicación de software que proporciona herramientas para escribir, compilar y depurar código. \\ \hline
IoT   & Internet of Things — Internet de las Cosas. Paradigma tecnológico que conecta dispositivos físicos a internet para la captura y transmisión automática de datos. \\ \hline
IPN   & Instituto Politécnico Nacional. Institución de educación superior pública de México. \\ \hline
JSON  & JavaScript Object Notation — Notación de Objetos de JavaScript. Formato ligero de intercambio de datos estructurados basado en pares clave-valor. \\ \hline
JWT   & JSON Web Token — Token Web en formato JSON. Estándar para la transmisión segura de información entre partes mediante tokens firmados digitalmente, utilizado en autenticación. \\ \hline
MUI   & Material-UI. Biblioteca de componentes visuales para React basada en el estándar de diseño Material Design de Google. \\ \hline
MXN   & Peso Mexicano. Código ISO 4217 de la moneda oficial de México. \\ \hline
PEPS  & Primeras Entradas, Primeras Salidas. Método de valuación de inventarios en el que los primeros artículos en entrar son los primeros en salir (equivalente al inglés FIFO). \\ \hline
PyME  & Pequeña y Mediana Empresa. Clasificación empresarial según el número de empleados y nivel de ingresos. \\ \hline
REST  & Representational State Transfer — Transferencia de Estado Representacional. Estilo de arquitectura de software para el diseño de servicios web mediante el protocolo HTTP. \\ \hline
RFID  & Radio Frequency Identification — Identificación por Radiofrecuencia. Tecnología que permite identificar objetos de forma automática mediante ondas de radio. \\ \hline
SPI   & Serial Peripheral Interface — Interfaz Periférica Serial. Protocolo de comunicación síncrona utilizado para la transferencia de datos entre microcontroladores y periféricos. \\ \hline
SQL   & Structured Query Language — Lenguaje de Consulta Estructurado. Lenguaje estándar para la gestión y manipulación de bases de datos relacionales. \\ \hline
USB   & Universal Serial Bus — Bus Serie Universal. Estándar de conexión que permite la comunicación y alimentación eléctrica entre dispositivos electrónicos. \\ \hline
XP    & Extreme Programming — Programación Extrema. Metodología ágil de desarrollo de software que enfatiza ciclos cortos de desarrollo y participación activa del usuario. \\ \hline

\end{longtable}
